\begin{enpars}

This book is a welcome resource to facilitate the chant of Vespers called for in parishes on Sundays and major feasts by the Second Vatican Council, so that all may pray more perfectly this evening sacrifice more perfectly, both as to their interior recollection as to the external celebration (\textit{Sacrosanctum Concilium} 100; 99), according to the Roman office published in 1911 by Pope Saint Pius \textsc{x}. The Sundays of the year, the major feasts of the temporal cycle and those fixed to Sunday, and those feasts which can take the place of the Sunday in occurrence or in concurrence are included.

%The feasts of the Queenship of the Blessed Virgin Mary and of the Immaculate Heart of Mary are included in an appendix.

%The feasts of the Queenship of the Blessed Virgin Mary and of the Immaculate Heart of Mary are included in an appendix, as is the revised office of the Assumption.

The commons needed to pray the office of Vespers on every Sunday are thus found in this volume, including the common of the Blessed Virgin Mary, as her status as patronal saint includes several offices which in the Roman rite take their parts mostly from the common in addition to the commons of Apostles and Evangelists both outside and in Paschal Time; although the Office of the Dead is prayed in choir on Sunday only when November 2 falls on this day, the usual collects are provided in addition to those prayed at the office following Vespers of All Saints, or those of Sunday when the Commemoration of All the Faithful Departed is transferred to November 3. The office of the dedication is also included, as there are feasts of the dedication of a church on the universal calendar, taking the place of the Sunday as feasts of the Lord.

The chief sources are the \textit{Liber Usualis} edited by the monks of Saint-Pierre de Solesmes (France) and the \textit{Liber antiphonarius} with the additions of the same monks. The syllabification follows the pattern used in these books, insofar as it can be determined. In order to speed the publication and dissemination of this book, the rhythmic signs are kept; since these remain in widespread use, their omission seems more problematic than their retention.

The psalms are marked with the major accents and preparatory syllables needed to chant in the various tones; these syllables are in italic and bold type respectively. The diaresis is not printed for aesthetic reasons, as it is unnecessary. In words of Hebrew origin such as ``Israel'' it is not possible for a diphthong to occur; note that in ecclesiastical Latin, ``Aaron'' has three syllables. The psalms are listed as frequently as possible to make the book as easy to use as possible without excluding the possibility of a page turn where the psalm is printed in the same tone, with the same ending. Where the psalm tone is marked with a *, as in tones 4A* and 8G*, one may freely omit the second note of the podatus, singing only the final of the mode, the one marked by the capital letter, as if the tone without the star was assigned to the antiphon. For other texts, notably the hymn and versicle of the season, or of I or II Vespers when this is in common, as many cross-references as possible are made; for the sake of space, however, the Canticle of the Blessed Virgin Mary is given only by tone with the appropriate endings as in the \textit{Liber Usualis.}

The tones required by the celebrant are not given; only the metrum and flex are indicated as in the chant editions, since the tone of the chapter never changes whereas one has a choice at the collect, although it is to be noted that the \textit{tonus solemnis,} from the \textit{toni antiqui ad libitum} in the common tones of the Vatican Edition, is more frequently used, except in the Office of the Dead and during the Triduum, when the oration is sung \textit{recto tono.} The ferial tones for the oration sung after the Marian antiphons are also found in the liturgical books.

Any commemoration of a Sunday or of an octave uses the versicle of the Sunday or of the feast but in the simple tone, with the ending formula of a descending minor third; dactyls such as ``Dóminus'' are customarily sung witth the inflection made on the second and third syllables. For Hebrew words and monosyllables, the usual formula is most commonly followed.

%%maybe to change with respect to hymn tones but hey, it's OUR book

The choices conform to the customs of the Church of the Assumption (Nashville, Tennessee, United States) influenced by the usage of the Institute of Christ the King Sovereign Priest and of other French churches using the classical Roman office. The tones of the hymns are given in full, but there is no obligation to use the tone listed first in the liturgical books. Since the usage of the solemn tone of the Gospel canticle and the \textit{Benedicamus Domino} tones are not absolutely defined, nor are the assignments free of any gaps, cross-references to a particular tone are avoided in order to increase the possibilities for using this book. But the following customs may be worth following:

\begin{itemize}
 \item The solemn tone at the beginning of Vespers is used on doubles of the \textsc{i} class, and the festal tone is used on doubles of the \textsc{ii} class and on feasts of the rank of dobule major, in addition to all other Sundays.
 \item The solemn tone of the \textit{Ma\-gnificat} is used on doubles of both the \textsc{i} and \textsc{ii} class. It is always used on the days of Advent where an ``O'' antiphon is sung.
 \item  At \textsc{ii} Vespers, the formula for the \textit{Benedicamus Domino} in mode \textsc{vi} is reserved for doubles of the \textsc{ii} class and less important feasts of our Lord, whereas the solemn tone in mode \textsc{v} is sung on the more important doubles of the \textsc{i} class. 
 \item
 The Marian tone in mode \textsc{i} is sung for all feasts of the Blessed Virgin Mary, even though the antiphonal calls for the tone of solemn feasts in certain cases (at a minimum, August 15 and December 8)
 \item
 The tone in mode \textsc{ii} for ``days within octaves'' is used during all octaves outside of Paschal Time, when the office is of the Sunday within the octave or of the octave day, and on the feast of the Holy Family.
 \item
 The solemn tone of the versicle after the hymn is used on feasts of the universal church where it is printed in the antiphonal; otherwise, the ordinary tone is always used.
 \item As to the tone of the hymn, the custom of the Church of the Assumption is to use the second tone of the hymn of Sunday Vespers and that in mode \textsc{iv} at Vespers of Paschal Time; the first tone is otherwise chosen, including the tone in mode \textsc{i} for feasts of the Blessed Virgin Mary no matter the rank, since in the Roman rite, one is free to choose any of the three on feasts other than on the simple feast of Our Lady \textit{in Sabbato.}
\end{itemize}

Some directions are given at Vespers of Sundays throughout the year. Those largely conforms to he usage of the Institute of Christ the King and to the rite described by Stercky in the \textit{Manuel de liturgie et cérémonial selon le rit romain}. However, directions for the cantors, the acolytes and the master of ceremonies, the assistants in cope, and the celebrant are omitted, as Vespers need not be sung solemnly or with acolytes, and there is wider discrepancy in custom in what concerns the actions of the cantors.

One may substitute ``cantor'' for ``cantor, and the schola'' or ``one side'' or ``men'' (or ``women'') and vice-versa with respect to the congregation, depending on whether they are invited to sing with the entire schola, or divided by sex, or by side, regardless of sex. Naturally, one may employ more cantors depending on the solemnity required and the number of cantors at one's disposal.

Otherwise, general knowledge of the rubrics is required; the chapter will always conclude with ``Deo grátias,'' and the collect is always preceded by ``Orémus'' unless otherwise noted. These are omitted along with the full conclusions in translation; the celebrant must memorize the full conclusions in Latin. Rubrics as found in the antiphonal, governing prescribed commemorations in perpetual occurrence or in concurrence, the celebration of the office during an octave, and so on are found in the appropriate places.

Additionally, recourse to an ordo is necessary to determine what office is to be said, and in what manner, as well the commemorations or suffrages must be said, if any are indeed required. The suffrages and commemorations of a Sunday or of an octave  can be made with this book, but one will need to refer to the \textit{Liber antiphonarius} in order to make all of the commemorations which arise from the office of a saint or of the Advent or Lenten feria merely commemorated. The offices of Sundays are of semidouble rite, except when noted as double feasts. On feasts, the antiphon is printed once, before the psalm, in those cases in order to minimize repetitions and to economize space.
% In Advent, the antiphons are printed as on doubles for the same reason, but the office remains of semidouble rite.
%%to change semidouble part for Advent??
%%to change, perhaps: commemoration of feria is kind of essential

Some feasts are included due to their importance even though they never occur on Sunday or do so only in concurrence, including when transferred to the following day (or to a free day following the octaves of Easter and of Pentecost, according to the rubrics) as in the last week of Advent, in Lent or in Passiontide. The feasts of Corpus Christi and of the Sacred Heart are of such importance, and the following Sundays within the octave take the antiphons, psalms, and hymn from the feast. The solemnity of Saint Joseph in Paschaltide is similarly included, as the Guardian of the Redeemer is Patron of the Universal Church. In all three cases, this greatly simplifies the cross-references required to make the commemoration of the octave, but it is to be reiterated that one must refer to an adequate ordo.

%%check octave rubric
Generally speaking, these feasts are to be observed as doubles of the \textsc{i} class subject to the usual rubrics governing occurrence and concurrence, all with a common octave:
\begin{itemize}
\item Titular feast
\item Dedication of the church
\item Dedication of the cathedral
\item Principal patron or patrons equally patron of a place, of a diocese, of a province, or of a nation; the feast of the principal patron, or if there be several equally patron, is celebrated as a double of the first class with a common octave.
\end{itemize}

In this volume, one will find the common of the dedication and the offices which may take the place of a Sunday on the universal calendar but may be kept with even greater solemnity as the titular or patronal feast.

NB: Ordinarily, all clerics are bound to celebrate the patronal feast of a place and of the nation, but the patronal feast of the diocese substitutes for that of the place if one is not named; that of the episcopal city if the diocese lacks a patron. Secular clerics and religious following the diocesan calendar keep this feast with a common octave. The national patronal feast is obligatory in all cases, and that of the diocese is obligatory if it was formerly a holy day of obligation.

The default titular feast for churches dedicated to the Lord is the Transfiguration, while the titular or patronal feast is that of the Assumption, in the case of the Blessed Virgin Mary, and the primary feast if a saint has more than one feast day, unless otherwise specified in all cases.

%%common commemorations require fixing names of doctors (again)
%%Lenten feria not possible in this volume
%%Triduum appendix??

%%should we cut down Orémus/Deo gratias/translation of Let us pray and full conclusion? Those take up a lot of space!
\end{enpars}